\Needspace{10\baselineskip}
\pdfbookmark[0]{Notices. Opérations géodésiques et astronomiques. 1830.}{anzeigen-operations-1830}
\begin{center}
\rule{.85\textwidth}{.3pt}\par
Göttingische gelehrte Anzeigen.\quad 27 February 1830.\par
\rule{.85\textwidth}{.3pt}
\end{center}
\hspace*{\parindent}\textit{Opérations géodésiques et astronomiques pour la mesure d’un arc du parallèle moyen, exécutées en Piémont et en Savoie par une commission composée d’officiers de l’état major général et d’astronomes Piémontais et Autrichiens en} 1821, 1822, 1823. \textit{Milan, de l’imprimerie impériale et royale: Tome premier} 1825. 238 pages. \textit{Tome second} 1827. 412 pages, in quarto. With a separate booklet containing figures, maps and six panoramic views.

The idea for the great survey of an arc of longitude, of which the operations published in this work form a principal part, originally came from \textsc{Laplace}. Since 1802, French engineers had carried out extensive triangulation in northern Italy, initially for topographical and military
